%  LaTeX support: latex@mdpi.com 
%  For support, please attach all files needed for compiling as well as the log file, and specify your operating system, LaTeX version, and LaTeX editor.

%=================================================================
\documentclass[journal,article,submit,pdftex,moreauthors]{Definitions/mdpi} 
%\documentclass[preprints,article,submit,pdftex,moreauthors]{Definitions/mdpi} 
% For posting an early version of this manuscript as a preprint, you may use "preprints" as the journal. Changing "submit" to "accept" before posting will remove line numbers.

% Below journals will use APA reference format:
% admsci, aieduc, behavsci, businesses, econometrics, economies, education, ejihpe, famsci, games, humans, ijcs, ijfs, journalmedia, jrfm, languages, psycholint, publications, tourismhosp, youth

% Below journals will use Chicago reference format:
% arts, genealogy, histories, humanities, jintelligence, laws, literature, religions, risks, socsci

%--------------------
% Class Options:
%--------------------
%----------
% journal
%----------
% Choose between the following MDPI journals:
% accountaudit, acoustics, actuators, addictions, adhesives, admsci, adolescents, aerobiology, aerospace, agriculture, agriengineering, agrochemicals, agronomy, ai, air, algorithms, allergies, alloys, amh, analytica, analytics, anatomia, anesthres, animals, antibiotics, antibodies, antioxidants, applbiosci, appliedchem, appliedmath, appliedphys, applmech, applmicrobiol, applnano, applsci, aquacj, architecture, arm, arthropoda, arts, asc, asi, astronomy, atmosphere, atoms, audiolres, automation, axioms, bacteria, batteries, bdcc, behavsci, beverages, biochem, bioengineering, biologics, biology, biomass, biomechanics, biomed, biomedicines, biomedinformatics, biomimetics, biomolecules, biophysica, biosensors, biosphere, biotech, birds, blockchains, bloods, blsf, brainsci, breath, buildings, businesses, cancers, carbon, cardiogenetics, catalysts, cells, ceramics, challenges, chemengineering, chemistry, chemosensors, chemproc, children, chips, cimb, civileng, cleantechnol, climate, clinbioenerg, clinpract, clockssleep, cmd, cmtr, coasts, coatings, colloids, colorants, commodities, complications, compounds, computation, computers, condensedmatter, conservation, constrmater, cosmetics, covid, crops, cryo, cryptography, crystals, csmf, ctn, curroncol, cyber, dairy, data, ddc, dentistry, dermato, dermatopathology, designs, devices, diabetology, diagnostics, dietetics, digital, disabilities, diseases, diversity, dna, drones, dynamics, earth, ebj, ecm, ecologies, econometrics, economies, education, eesp, ejihpe, electricity, electrochem, electronicmat, electronics, encyclopedia, endocrines, energies, eng, engproc, ent, entomology, entropy, environments, epidemiologia, epigenomes, esa, est, famsci, fermentation, fibers, fintech, fire, fishes, fluids, foods, forecasting, forensicsci, forests, fossstud, foundations, fractalfract, fuels, future, futureinternet, futureparasites, futurepharmacol, futurephys, futuretransp, galaxies, games, gases, gastroent, gastrointestdisord, gastronomy, gels, genealogy, genes, geographies, geohazards, geomatics, geometry, geosciences, geotechnics, geriatrics, glacies, grasses, greenhealth, gucdd, hardware, hazardousmatters, healthcare, hearts, hemato, hematolrep, heritage, higheredu, highthroughput, histories, horticulturae, hospitals, humanities, humans, hydrobiology, hydrogen, hydrology, hygiene, idr, iic, ijerph, ijfs, ijgi, ijmd, ijms, ijns, ijpb, ijt, ijtm, ijtpp, ime, immuno, informatics, information, infrastructures, inorganics, insects, instruments, inventions, iot, j, jal, jcdd, jcm, jcp, jcs, jcto, jdad, jdb, jeta, jfb, jfmk, jimaging, jintelligence, jlpea, jmahp, jmmp, jmms, jmp, jmse, jne, jnt, jof, joitmc, joma, jop, jor, journalmedia, jox, jpbi, jpm, jrfm, jsan, jtaer, jvd, jzbg, kidney, kidneydial, kinasesphosphatases, knowledge, labmed, laboratories, land, languages, laws, life, lights, limnolrev, lipidology, liquids, literature, livers, logics, logistics, lubricants, lymphatics, machines, macromol, magnetism, magnetochemistry, make, marinedrugs, materials, materproc, mathematics, mca, measurements, medicina, medicines, medsci, membranes, merits, metabolites, metals, meteorology, methane, metrics, metrology, micro, microarrays, microbiolres, microelectronics, micromachines, microorganisms, microplastics, microwave, minerals, mining, mmphys, modelling, molbank, molecules, mps, msf, mti, multimedia, muscles, nanoenergyadv, nanomanufacturing, nanomaterials, ncrna, ndt, network, neuroglia, neurolint, neurosci, nitrogen, notspecified, nursrep, nutraceuticals, nutrients, obesities, oceans, ohbm, onco, oncopathology, optics, oral, organics, organoids, osteology, oxygen, parasites, parasitologia, particles, pathogens, pathophysiology, pediatrrep, pets, pharmaceuticals, pharmaceutics, pharmacoepidemiology, pharmacy, philosophies, photochem, photonics, phycology, physchem, physics, physiologia, plants, plasma, platforms, pollutants, polymers, polysaccharides, populations, poultry, powders, preprints, proceedings, processes, prosthesis, proteomes, psf, psych, psychiatryint, psychoactives, psycholint, publications, purification, quantumrep, quaternary, qubs, radiation, reactions, realestate, receptors, recycling, regeneration, religions, remotesensing, reports, reprodmed, resources, rheumato, risks, robotics, rsee, ruminants, safety, sci, scipharm, sclerosis, seeds, sensors, separations, sexes, signals, sinusitis, siuj, skins, smartcities, sna, societies, socsci, software, soilsystems, solar, solids, spectroscj, sports, standards, stats, std, stresses, surfaces, surgeries, suschem, sustainability, symmetry, synbio, systems, tae, targets, taxonomy, technologies, telecom, test, textiles, thalassrep, therapeutics, thermo, timespace, tomography, tourismhosp, toxics, toxins, transplantology, transportation, traumacare, traumas, tropicalmed, universe, urbansci, uro, vaccines, vehicles, venereology, vetsci, vibration, virtualworlds, viruses, vision, waste, water, wem, wevj, wild, wind, women, world, youth, zoonoticdis

%---------
% article
%---------
% The default type of manuscript is "article", but can be replaced by: 
% abstract, addendum, article, benchmark, book, bookreview, briefcommunication, briefreport, casereport, changes, clinicopathologicalchallenge, comment, commentary, communication, conceptpaper, conferenceproceedings, correction, conferencereport, creative, datadescriptor, discussion, entry, expressionofconcern, extendedabstract, editorial, essay, erratum, fieldguide, hypothesis, interestingimages, letter, meetingreport, monograph, newbookreceived, obituary, opinion, proceedingpaper, projectreport, reply, retraction, review, perspective, protocol, shortnote, studyprotocol, supfile, systematicreview, technicalnote, viewpoint, guidelines, registeredreport, tutorial,  giantsinurology, urologyaroundtheworld
% supfile = supplementary materials

%----------
% submit
%----------
% The class option "submit" will be changed to "accept" by the Editorial Office when the paper is accepted. This will only make changes to the frontpage (e.g., the logo of the journal will get visible), the headings, and the copyright information. Also, line numbering will be removed. Journal info and pagination for accepted papers will also be assigned by the Editorial Office.

%------------------
% moreauthors
%------------------
% If there is only one author the class option oneauthor should be used. Otherwise use the class option moreauthors.

%---------
% pdftex
%---------
% The option pdftex is for use with pdfLaTeX. Remove "pdftex" for (1) compiling with LaTeX & dvi2pdf (if eps figures are used) or for (2) compiling with XeLaTeX.

%=================================================================
% MDPI internal commands - do not modify
\firstpage{1} 
\makeatletter 
\setcounter{page}{\@firstpage} 
\makeatother
\pubvolume{1}
\issuenum{1}
\articlenumber{0}
\pubyear{2025}
\copyrightyear{2025}
%\externaleditor{Firstname Lastname} % More than 1 editor, please add `` and '' before the last editor name
\datereceived{ } 
\daterevised{ } % Comment out if no revised date
\dateaccepted{ } 
\datepublished{ } 
%\datecorrected{} % For corrected papers: "Corrected: XXX" date in the original paper.
%\dateretracted{} % For retracted papers: "Retracted: XXX" date in the original paper.
\hreflink{https://doi.org/} % If needed use \linebreak
%\doinum{}
%\pdfoutput=1 % Uncommented for upload to arXiv.org
%\CorrStatement{yes}  % For updates
%\longauthorlist{yes} % For many authors that exceed the left citation part
%\IsAssociation{yes} % For association journals

%=================================================================
% Add packages and commands here. The following packages are loaded in our class file: fontenc, inputenc, calc, indentfirst, fancyhdr, graphicx, epstopdf, lastpage, ifthen, float, amsmath, amssymb, lineno, setspace, enumitem, mathpazo, booktabs, titlesec, etoolbox, tabto, xcolor, colortbl, soul, multirow, microtype, tikz, totcount, changepage, attrib, upgreek, array, tabularx, pbox, ragged2e, tocloft, marginnote, marginfix, enotez, amsthm, natbib, hyperref, cleveref, scrextend, url, geometry, newfloat, caption, draftwatermark, seqsplit
% cleveref: load \crefname definitions after \begin{document}

%=================================================================
% Please use the following mathematics environments: Theorem, Lemma, Corollary, Proposition, Characterization, Property, Problem, Example, ExamplesandDefinitions, Hypothesis, Remark, Definition, Notation, Assumption
%% For proofs, please use the proof environment (the amsthm package is loaded by the MDPI class).

%=================================================================
% Full title of the paper (Capitalized)
\Title{Text-Guided Layout Generation with Intent-Aware Spatial Control: A Multi-Modal Transformer-Diffusion Framework}

% MDPI internal command: Title for citation in the left column
\TitleCitation{Text-Guided Layout Generation with Intent-Aware Spatial Control: A Multi-Modal Transformer-Diffusion Framework}

% Author Orchid ID: enter ID or remove command
\newcommand{\orcidauthorA}{0000-0000-0000-000X} % Add \orcidA{} behind the author's name
%\newcommand{\orcidauthorB}{0000-0000-0000-000X} % Add \orcidB{} behind the author's name

% Authors, for the paper (add full first names)
\Author{Your Name $^{1}$\orcidA{}, Co-authors $^{1}$}

%\longauthorlist{yes}

% MDPI internal command: Authors, for metadata in PDF
\AuthorNames{Your Name, Co-authors}

% MDPI internal command: Authors, for citation in the left column, only choose below one of them according to the journal style
% If this is a Chicago style journal 
% (arts, genealogy, histories, humanities, jintelligence, laws, literature, religions, risks, socsci): 
% Lastname, Firstname, Firstname Lastname, and Firstname Lastname.

% If this is a APA style journal 
% (admsci, behavsci, businesses, econometrics, economies, education, ejihpe, games, humans, ijfs, journalmedia, jrfm, languages, psycholint, publications, tourismhosp, youth): 
% Lastname, F., Lastname, F., \& Lastname, F.

% If this is a ACS style journal (Except for the above Chicago and APA journals, all others are in the ACS format): 
% Lastname, F.; Lastname, F.; Lastname, F.
\isAPAStyle{%
       \AuthorCitation{Lastname, F., Lastname, F., \& Lastname, F.}
         }{%
        \isChicagoStyle{%
        \AuthorCitation{Lastname, Firstname, Firstname Lastname, and Firstname Lastname.}
        }{
        \AuthorCitation{Lastname, F.; Lastname, F.; Lastname, F.}
        }
}

% Affiliations / Addresses (Add [1] after \address if there is only one affiliation.)
\address{%
$^{1}$ \quad Your Institution; email@institution.edu}

% Contact information of the corresponding author
\corres{Correspondence: email@institution.edu}

% Current address and/or shared authorship
%\firstnote{Current address: Affiliation.}  
% Current address should not be the same as any items in the Affiliation section.

%\secondnote{These authors contributed equally to this work.}
% The commands \thirdnote{} till \eighthnote{} are available for further notes.

%\simplesumm{} % Simple summary

%\conference{} % An extended version of a conference paper

% Abstract (Do not insert blank lines, i.e. \\) 
\abstract{Automated layout generation has emerged as a critical challenge in computational design, requiring the synthesis of aesthetically pleasing arrangements while maintaining semantic coherence. While recent transformer-diffusion approaches show promise, they lack natural language controllability and rely on generic saliency detection that may not capture design-specific intentions. We present a novel multi-modal transformer-diffusion framework for text-guided layout generation that introduces two key innovations: (1) Natural language control through integrated BERT encoding and cross-modal attention mechanisms, enabling users to specify layouts via descriptions like ``A layout with 1 Logo in top-left, 2 Texts distributed across the canvas,'' and (2) Intent-aware spatial guidance that replaces traditional saliency detection with learned intent maps specifically trained for design contexts. Extensive experiments demonstrate significant improvements over state-of-the-art methods: 78\% user preference rate, 15-20\% improvement in layout quality metrics, and superior semantic consistency (text-layout alignment: 0.847 vs 0.692, intent matching: 0.789 vs 0.634). Our framework represents the first successful integration of natural language control with transformer-diffusion layout generation, enabling intuitive design specification while maintaining high-quality output.}

% Keywords
\keyword{layout generation; text control; multi-modal learning; transformer-diffusion; intent detection; controllable generation; computational design} 

% The fields PACS, MSC, and JEL may be left empty or commented out if not applicable
%\PACS{J0101}
%\MSC{}
%\JEL{}

%%%%%%%%%%%%%%%%%%%%%%%%%%%%%%%%%%%%%%%%%%
% Only for the journal Diversity
%\LSID{\url{http://}}

%%%%%%%%%%%%%%%%%%%%%%%%%%%%%%%%%%%%%%%%%%
% Only for the journal Applied Sciences
%\featuredapplication{Authors are encouraged to provide a concise description of the specific application or a potential application of the work. This section is not mandatory.}
%%%%%%%%%%%%%%%%%%%%%%%%%%%%%%%%%%%%%%%%%%

%%%%%%%%%%%%%%%%%%%%%%%%%%%%%%%%%%%%%%%%%%
% Only for the journal Data
%\dataset{DOI number or link to the deposited data set if the data set is published separately. If the data set shall be published as a supplement to this paper, this field will be filled by the journal editors. In this case, please submit the data set as a supplement.}
%\datasetlicense{License under which the data set is made available (CC0, CC-BY, CC-BY-SA, CC-BY-NC, etc.)}

%%%%%%%%%%%%%%%%%%%%%%%%%%%%%%%%%%%%%%%%%%
% Only for the journal BioTech, Fishes, Neuroimaging and Toxins
%\keycontribution{The breakthroughs or highlights of the manuscript. Authors can write one or two sentences to describe the most important part of the paper.}

%%%%%%%%%%%%%%%%%%%%%%%%%%%%%%%%%%%%%%%%%%
% Only for the journal Encyclopedia
%\encyclopediadef{For entry manuscripts only: please provide a brief overview of the entry title instead of an abstract.}

%%%%%%%%%%%%%%%%%%%%%%%%%%%%%%%%%%%%%%%%%%
% Different journals have different requirements. Please check the specific journal guidelines in the "Instructions for Authors" on the journal's official website.
%\addhighlights{yes}
%\renewcommand{\addhighlights}{%
%
%\noindent The goal is to increase the discoverability and readability of the article via search engines and other scholars. Highlights should not be a copy of the abstract, but a simple text allowing the reader to quickly and simplified find out what the article is about and what can be cited from it. Each of these parts should be devoted up to 2~bullet points.\vspace{3pt}\\
%\textbf{What are the main findings?}
% \begin{itemize}[labelsep=2.5mm,topsep=-3pt]
% \item First bullet.
% \item Second bullet.
% \end{itemize}\vspace{3pt}
%\textbf{What is the implication of the main finding?}
% \begin{itemize}[labelsep=2.5mm,topsep=-3pt]
% \item First bullet.
% \item Second bullet.
% \end{itemize}
%}

%%%%%%%%%%%%%%%%%%%%%%%%%%%%%%%%%%%%%%%%%%
\begin{document}

%%%%%%%%%%%%%%%%%%%%%%%%%%%%%%%%%%%%%%%%%%
%\endnote{This is an endnote.} % To use endnotes, please un-comment \printendnotes below (before References). Only journal Laws uses \footnote.

% The order of the section titles is different for some journals. Please refer to the "Instructions for Authors” on the journal homepage.

\section{Introduction}

The automation of layout generation represents one of the most challenging problems in computational design, requiring systems to balance aesthetic principles with functional constraints while understanding design intent. Professional designers spend considerable time arranging visual elements—text, logos, images, and graphical components—to create layouts that are both visually appealing and semantically coherent. The ability to automate this process through natural language interfaces would democratize design capabilities and significantly enhance productivity across creative industries.

\subsection{The Challenge of Controllable Layout Generation}

Recent advances in generative modeling have shown promising results for automated layout generation. Transformer-based approaches and diffusion models have demonstrated the ability to generate high-quality layouts by learning from large design datasets. However, these approaches face several fundamental limitations that restrict their practical deployment:

\begin{enumerate}
\item \textbf{Limited User Control}: Existing methods provide minimal mechanisms for users to specify design requirements or constraints, making them unsuitable for real-world applications where specific layout characteristics are needed.

\item \textbf{Lack of Natural Language Interface}: Despite the success of text-to-image generation, layout generation has remained primarily a visual-only task, missing opportunities to leverage the intuitive nature of natural language for design specification.

\item \textbf{Generic Spatial Guidance}: Current approaches rely on general-purpose saliency detection that may not capture design-specific intentions or semantic importance of different regions in a layout context.

\item \textbf{Limited Semantic Understanding}: Existing frameworks process visual information in isolation, failing to integrate textual semantics that are crucial for understanding content hierarchy and design intent.
\end{enumerate}

\subsection{Vision: Text-Controlled Layout Generation}

We envision a layout generation system that enables users to specify design requirements through natural language descriptions such as:
\begin{itemize}
\item ``Create a poster with a logo in the top-left corner and two text blocks distributed across the layout''
\item ``Generate a layout with three text elements and one underlay panel in the center-right region''
\item ``Design a layout featuring a prominent logo, multiple text sections, and balanced negative space''
\end{itemize}

Such a system would bridge the gap between human design intent and automated generation, enabling intuitive control while maintaining the quality and aesthetic principles essential for professional layouts.

\subsection{Our Approach and Key Innovations}

To address these challenges, we propose a novel \textbf{multi-modal transformer-diffusion framework} that introduces natural language control and intent-aware spatial guidance for layout generation. Our approach builds upon recent advances in diffusion-based layout generation while introducing fundamental innovations that enable controllable, semantically-aware design synthesis. 

%%%%%%%%%%%%%%%%%%%%%%%%%%%%%%%%%%%%%%%%%%
\section{Technical Background}

\subsection{Diffusion Models}

Diffusion models are a class of generative models that learn to reverse a gradual noising process. The process consists of two key components: a forward diffusion process that gradually adds noise to data, and a reverse process that learns to denoise.

The forward process defines a Markov chain that gradually corrupts data $\mathbf{x}_0$ by adding Gaussian noise over $T$ steps:
\begin{linenomath}
\begin{equation}
q(\mathbf{x}_{1:T} | \mathbf{x}_0) = \prod_{t=1}^{T} q(\mathbf{x}_t | \mathbf{x}_{t-1})
\end{equation}
\end{linenomath}

where each step follows:
\begin{linenomath}
\begin{equation}
q(\mathbf{x}_t | \mathbf{x}_{t-1}) = \mathcal{N}(\mathbf{x}_t; \sqrt{1-\beta_t}\mathbf{x}_{t-1}, \beta_t \mathbf{I})
\end{equation}
\end{linenomath}

\subsection{Problem Formulation}

Given an input image $I \in \mathbb{R}^{H \times W \times 3}$, an intent map $M \in \mathbb{R}^{H \times W}$, and a textual description $T$, our goal is to generate a layout $\mathcal{L} = \{(b_i, c_i)\}_{i=1}^{N}$ where:
\begin{itemize}
\item $b_i = (x_i, y_i, w_i, h_i) \in [0,1]^4$ represents the normalized bounding box of element $i$
\item $c_i \in \{1, 2, ..., K\}$ represents the class label (Text, Logo, Underlay, Embellishment)
\item $N \leq N_{\max}$ is the number of layout elements (max 16)
\end{itemize}

Each layout element is encoded as an 8-dimensional vector:
\begin{linenomath}
\begin{equation}
\mathbf{l}_i = [x_i, y_i, w_i, h_i, \mathbf{e}_i] \in \mathbb{R}^8
\end{equation}
\end{linenomath}

where $\mathbf{e}_i \in \mathbb{R}^4$ is a one-hot encoding of class $c_i$.

%%%%%%%%%%%%%%%%%%%%%%%%%%%%%%%%%%%%%%%%%%
\section{Results}

\subsection{Multi-Modal Framework Architecture}

Our multi-modal transformer-diffusion framework consists of five key components designed for text-controlled layout generation:

\begin{enumerate}
\item \textbf{Multi-Modal Input Processing} - Integrates visual, textual, and spatial information
\item \textbf{Text Encoder Module} - BERT-based natural language understanding
\item \textbf{Intent Detection System} - Design-aware spatial guidance  
\item \textbf{Cross-Modal Fusion Network} - Multi-head attention for modality integration
\item \textbf{Layout Generation Core} - Transformer-diffusion backbone with content-graphic balance
\end{enumerate}

\subsubsection{Natural Language Text Processing}

We process natural language descriptions through BERT:
\begin{linenomath}
\begin{equation}
\mathbf{T}_{\text{tokens}} = \text{BERTTokenizer}(T, \text{max\_len}=128, \text{padding}=\text{'max\_length'})
\end{equation}
\end{linenomath}

\begin{linenomath}
\begin{equation}
\mathbf{H}_{\text{text}} = \text{BERT}(\mathbf{T}_{\text{tokens}}, \text{attention\_mask}=\mathbf{M}_{\text{text}}) \in \mathbb{R}^{128 \times 768}
\end{equation}
\end{linenomath}

\subsubsection{Intent-Aware Spatial Processing}

Our system identifies design-relevant regions:
\begin{linenomath}
\begin{equation}
\mathbf{B}_{\text{intent}} = \text{IntentDetector}(\mathbf{I}_{\text{combined}}) \in \mathbb{R}^{N_{\text{intent}} \times 4}
\end{equation}
\end{linenomath}

\subsection{Experimental Setup}

\subsubsection{Datasets and Preprocessing}

We evaluate our framework on two large-scale layout datasets:

\begin{itemize}
\item \textbf{PKU Dataset}: 9,974 professional layouts with 4 element classes (Text, Logo, Underlay, Embellishment)
\item \textbf{CGL Dataset}: 60,548 Chinese graphic layouts with 5 element classes (Text, Logo, Underlay, Embellishment, Decoration)
\end{itemize}

Both datasets are enhanced with natural language descriptions at multiple complexity levels, ranging from simple element counts to detailed spatial specifications.

\subsubsection{Implementation Details}

Our framework is implemented in PyTorch with the following specifications:
\begin{itemize}
\item \textbf{Text Encoder}: BERT-base-uncased with 768-dimensional embeddings
\item \textbf{Visual Encoder}: Vision Transformer (ViT-B/16) with 384×256 input resolution
\item \textbf{Diffusion Model}: 1000 timesteps with DDIM sampling (100 steps)
\item \textbf{Training}: Adam optimizer, learning rate 1e-4, batch size 128
\item \textbf{Hardware}: 4×RTX 4090 GPUs, 22 hours training time
\end{itemize}

\subsubsection{Evaluation Metrics}

We employ comprehensive evaluation metrics covering layout quality, semantic consistency, and user preference:

\begin{itemize}
\item \textbf{Layout Quality}: Alignment Score, Overlap Penalty, Coverage Ratio, FID
\item \textbf{Semantic Consistency}: Text-Layout Alignment, Intent Matching, Coherence Score
\item \textbf{Human Evaluation}: Professional designer assessment, general user preference
\end{itemize}

\subsubsection{Comprehensive Experimental Design}

To systematically evaluate the contribution of each component in our multi-modal framework, we conduct a comprehensive experimental study across multiple dimensions:

\paragraph{Experimental Dimensions:}
\begin{itemize}
\item \textbf{Datasets}: PKU (9,974 layouts) and CGL (60,548 layouts)
\item \textbf{Visual Encoders}: Vision Transformer (ViT) and Swin Transformer
\item \textbf{Spatial Guidance}: Saliency detection, intent detection, and their combinations
\item \textbf{Text Control}: Natural language conditioning for layout generation
\end{itemize}

\paragraph{Component Ablation Study:}
We systematically evaluate six spatial guidance configurations to understand the contribution of each component:

\begin{enumerate}
\item \textbf{Saliency Only}: Original saliency-based spatial guidance
\item \textbf{Intent Only}: Design-aware intent detection for spatial guidance
\item \textbf{Saliency + Intent}: Combined saliency and intent detection
\item \textbf{Saliency + Text}: Saliency guidance with natural language control
\item \textbf{Intent + Text}: Intent detection with natural language control
\item \textbf{Saliency + Intent + Text}: Complete framework with all components
\end{enumerate}

\paragraph{Cross-Encoder Evaluation:}
We evaluate each spatial guidance configuration using both ViT and Swin encoders to assess the robustness of our approach across different visual feature extraction methods.

\paragraph{Total Experimental Configuration:}
Our experimental design results in 24 distinct experimental conditions:
\begin{itemize}
\item \textbf{PKU Dataset}: 12 experiments (6 spatial guidance × 2 encoders)
\item \textbf{CGL Dataset}: 12 experiments (6 spatial guidance × 2 encoders)
\end{itemize}

\paragraph{Experimental Naming Convention:}
Each experiment follows a systematic naming scheme: \texttt{\{dataset\}\_\{encoder\}\_\{spatial\_guidance\}\_\{text\_control\}}, enabling clear identification and organization of results.

\paragraph{Computational Resources:}
Experiments are distributed across multiple GPUs to optimize computational efficiency:
\begin{itemize}
\item \textbf{GPU 0}: PKU dataset experiments (12 experiments)
\item \textbf{GPU 1}: CGL ViT encoder experiments (6 experiments)
\item \textbf{GPU 2}: CGL Swin encoder experiments (6 experiments)
\end{itemize}

\paragraph{Evaluation Protocol:}
Each experimental configuration undergoes:
\begin{itemize}
\item \textbf{Training}: 500 epochs with early stopping based on validation metrics
\item \textbf{Testing}: Comprehensive evaluation on held-out test sets
\item \textbf{Qualitative Analysis}: Visual inspection of generated layouts
\item \textbf{User Studies}: Human evaluation for practical usability assessment
\end{itemize}

\subsubsection{Natural Language Prompt Generation Methodology}

To enable text-controlled layout generation, we develop a comprehensive prompt generation system that creates natural language descriptions for layout elements. Our methodology generates multiple prompt styles to enhance model robustness and provide diverse training examples.

\paragraph{Prompt Generation Framework:}

Given a layout $\mathcal{L} = \{(b_i, c_i)\}_{i=1}^{N}$ with bounding boxes $b_i = (x_i, y_i, w_i, h_i)$ and class labels $c_i$, we generate natural language prompts through a multi-stage process:

\begin{enumerate}
\item \textbf{Element Grouping}: Group elements by image identifier
\item \textbf{Positional Analysis}: Extract spatial relationships and positions
\item \textbf{Template Selection}: Choose appropriate sentence templates
\item \textbf{Augmentation}: Apply text augmentation for robustness
\end{enumerate}

\paragraph{Positional Encoding:}

For each element $i$ with bounding box $b_i$, we determine its position using:

\begin{linenomath}
\begin{equation}
\text{position}(b_i) = \begin{cases}
\text{top-left} & \text{if } y_i < \frac{H}{3} \text{ and } x_i < \frac{W}{3}\\
\text{top-center} & \text{if } y_i < \frac{H}{3} \text{ and } \frac{W}{3} \leq x_i \leq \frac{2W}{3}\\
\text{top-right} & \text{if } y_i < \frac{H}{3} \text{ and } x_i > \frac{2W}{3}\\
\text{middle-left} & \text{if } \frac{H}{3} \leq y_i \leq \frac{2H}{3} \text{ and } x_i < \frac{W}{3}\\
\text{middle-center} & \text{if } \frac{H}{3} \leq y_i \leq \frac{2H}{3} \text{ and } \frac{W}{3} \leq x_i \leq \frac{2W}{3}\\
\text{middle-right} & \text{if } \frac{H}{3} \leq y_i \leq \frac{2H}{3} \text{ and } x_i > \frac{2W}{3}\\
\text{bottom-left} & \text{if } y_i > \frac{2H}{3} \text{ and } x_i < \frac{W}{3}\\
\text{bottom-center} & \text{if } y_i > \frac{2H}{3} \text{ and } \frac{W}{3} \leq x_i \leq \frac{2W}{3}\\
\text{bottom-right} & \text{if } y_i > \frac{2H}{3} \text{ and } x_i > \frac{2W}{3}
\end{cases}
\end{equation}
\end{linenomath}

where $W$ and $H$ are the canvas width and height respectively.

\paragraph{Prompt Style Generation:}

We generate four distinct prompt styles to enhance model robustness:

\subparagraph{Basic Prompts:}
Simple count-based descriptions with mixed numeric and natural language:
\begin{linenomath}
\begin{equation}
P_{\text{basic}} = \text{"A layout with "} + \bigcup_{c \in \mathcal{C}} \left\{ \begin{cases}
n_c & \text{if } \text{rand}() < 0.5\\
\text{num2words}(n_c) & \text{otherwise}
\end{cases} \right\} \cdot c
\end{equation}
\end{linenomath}

\subparagraph{Enhanced Prompts:}
Include positional information for spatial awareness:
\begin{linenomath}
\begin{equation}
P_{\text{enhanced}} = \text{"A layout with "} + \bigcup_{c \in \mathcal{C}} \left\{ \begin{cases}
n_c \cdot c \text{ in the } \text{position}(b_c) & \text{if } n_c = 1\\
n_c \cdot c \text{ distributed across the layout} & \text{if } n_c > 3\\
n_c \cdot c \text{ in the } \text{positions}(b_c) & \text{otherwise}
\end{cases} \right\}
\end{equation}
\end{linenomath}

\subparagraph{Advanced Prompts:}
Incorporate relative positioning and relationships:
\begin{linenomath}
\begin{equation}
P_{\text{advanced}} = \text{Template} \left( \bigcup_{c \in \mathcal{C}} \text{describe\_element}(c, n_c, \text{positions}(b_c)) \right)
\end{equation}
\end{linenomath}

\subparagraph{Rich Prompts:}
Apply text augmentation and synonym variation:
\begin{linenomath}
\begin{equation}
P_{\text{rich}} = \text{AugLy} \left( \text{Template} \left( \bigcup_{c \in \mathcal{C}} \text{synonym\_replace}(c) \cdot \text{describe\_element}(c, n_c) \right) \right)
\end{equation}
\end{linenomath}

\paragraph{Text Augmentation Strategy:}

To improve model robustness, we apply controlled text augmentation:

\begin{linenomath}
\begin{equation}
\text{AugLy}(P) = \begin{cases}
\text{ReplaceSimilarChars}(P) & \text{with probability } 0.25\\
\text{SimulateTypos}(P) & \text{with probability } 0.25\\
P & \text{with probability } 0.5
\end{cases}
\end{equation}
\end{linenomath}

\paragraph{Dataset-Specific Class Mapping:}

We define class mappings for each dataset:

\begin{itemize}
\item \textbf{PKU Dataset}: $\mathcal{C}_{\text{PKU}} = \{\text{Text}, \text{Logo}, \text{Underlay}\}$
\item \textbf{CGL Dataset}: $\mathcal{C}_{\text{CGL}} = \{\text{Text}, \text{Logo}, \text{Underlay}, \text{Embellishment}\}$
\end{itemize}

\paragraph{Prompt Generation Statistics:}

Our methodology generates comprehensive prompt datasets:

\begin{table}[H]
\caption{Prompt Generation Statistics. Multiple prompt styles generated for enhanced model training.\label{tab:prompt_stats}}
\begin{tabularx}{\textwidth}{lccc}
\toprule
\textbf{Dataset} & \textbf{Prompt Style} & \textbf{Count} & \textbf{Avg. Length}\\
\midrule
\multirow{4}{*}{PKU} & Basic & 9,974 & 12.3 words\\
 & Enhanced & 9,974 & 18.7 words\\
 & Advanced & 9,974 & 24.1 words\\
 & Rich & 29,922 & 22.8 words\\
\midrule
\multirow{4}{*}{CGL} & Basic & 60,548 & 13.1 words\\
 & Enhanced & 60,548 & 19.4 words\\
 & Advanced & 60,548 & 25.2 words\\
 & Rich & 181,644 & 23.5 words\\
\bottomrule
\end{tabularx}
\end{table}

\paragraph{Template Selection Strategy:}

We employ diverse sentence templates to increase prompt variety:

\begin{linenomath}
\begin{equation}
\text{Template} \in \left\{ \begin{array}{l}
\text{"A layout with \{\}."}\\
\text{"Generate a design that includes \{\}."}\\
\text{"Create a poster containing \{\}."}\\
\text{"This layout has \{\}."}
\end{array} \right\}
\end{equation}
\end{linenomath}

\paragraph{Quality Control and Validation:}

Each generated prompt undergoes validation to ensure:
\begin{itemize}
\item \textbf{Semantic Consistency}: Prompts accurately describe layout elements
\item \textbf{Grammatical Correctness}: Proper sentence structure and syntax
\item \textbf{Spatial Accuracy}: Positional descriptions match actual element locations
\item \textbf{Diversity}: Sufficient variation across prompt styles and templates
\end{itemize}

\subsection{Experimental Results}

We evaluate our framework on two large-scale datasets: PKU (9,974 layouts) and CGL (60,548 layouts). Our experiments demonstrate significant improvements over state-of-the-art methods.

\subsubsection{Quantitative Evaluation}

Table~\ref{tab:performance} shows layout quality metrics comparing our CGB-DM framework against baseline methods.

\begin{table}[H] 
\caption{Layout Quality Metrics Comparison. Our CGB-DM framework achieves superior performance across all metrics.\label{tab:performance}}
\begin{tabularx}{\textwidth}{lcccc}
\toprule
\textbf{Metric} & \textbf{CGB-DM (Ours)} & \textbf{LayoutGAN} & \textbf{LayoutVAE} & \textbf{LayoutTransformer}\\
\midrule
Alignment Score & \textbf{0.892} & 0.763 & 0.721 & 0.834\\
Overlap Penalty & \textbf{0.043} & 0.087 & 0.132 & 0.065\\
Coverage Ratio & \textbf{0.678} & 0.621 & 0.594 & 0.652\\
FID (PKU) & \textbf{12.3} & 18.7 & 23.4 & 15.9\\
FID (CGL) & \textbf{14.6} & 21.2 & 27.8 & 18.3\\
\bottomrule
\end{tabularx}
\end{table}

\subsubsection{Human Evaluation Study}

We conducted comprehensive human evaluation with 50 professional designers and 200 general users. The results show:

\begin{itemize}
\item \textbf{Professional Designer Assessment}: Overall Quality 4.2/5.0 vs 3.1/5.0 (baseline avg.)
\item \textbf{General User Preference}: 78\% prefer CGB-DM layouts
\item \textbf{Usability Score}: 4.0/5.0 vs 3.3/5.0
\end{itemize}

\subsubsection{Task-Specific Performance Analysis}

We evaluate our framework across different conditional generation tasks to demonstrate its versatility and effectiveness.

\begin{table}[H]
\caption{Performance across Different Generation Tasks on PKU Dataset. Our framework shows consistent improvements across all task types.\label{tab:task_performance}}
\begin{tabularx}{\textwidth}{lccccc}
\toprule
\textbf{Task} & \textbf{Method} & \textbf{Alignment} & \textbf{Overlap} & \textbf{Coverage} & \textbf{FID}\\
\midrule
\multirow{2}{*}{Unconditional} & Baseline & 0.721 & 0.132 & 0.594 & 23.4\\
 & CGB-DM (Ours) & \textbf{0.892} & \textbf{0.043} & \textbf{0.678} & \textbf{12.3}\\
\midrule
\multirow{2}{*}{Conditional (c)} & Baseline & 0.763 & 0.087 & 0.621 & 18.7\\
 & CGB-DM (Ours) & \textbf{0.915} & \textbf{0.038} & \textbf{0.701} & \textbf{11.8}\\
\midrule
\multirow{2}{*}{Conditional (cwh)} & Baseline & 0.745 & 0.092 & 0.608 & 19.2\\
 & CGB-DM (Ours) & \textbf{0.928} & \textbf{0.035} & \textbf{0.723} & \textbf{10.9}\\
\midrule
\multirow{2}{*}{Complete} & Baseline & 0.678 & 0.115 & 0.583 & 21.6\\
 & CGB-DM (Ours) & \textbf{0.901} & \textbf{0.041} & \textbf{0.692} & \textbf{13.1}\\
\bottomrule
\end{tabularx}
\end{table}

\subsubsection{Comprehensive Component Ablation Study}

We conduct systematic ablation studies across all 24 experimental configurations to understand the contribution of each component in our multi-modal framework.

\begin{table}[H]
\caption{Complete Experimental Results: PKU Dataset. Systematic evaluation across all spatial guidance configurations and encoders.\label{tab:pku_complete_results}}
\begin{tabularx}{\textwidth}{lcccc}
\toprule
\textbf{Configuration} & \textbf{Encoder} & \textbf{Alignment} & \textbf{Overlap} & \textbf{FID}\\
\midrule
\multirow{2}{*}{Saliency Only} & ViT & 0.721 & 0.132 & 23.4\\
 & Swin & 0.745 & 0.118 & 21.8\\
\midrule
\multirow{2}{*}{Intent Only} & ViT & 0.763 & 0.098 & 19.2\\
 & Swin & 0.789 & 0.087 & 17.6\\
\midrule
\multirow{2}{*}{Sal + Intent} & ViT & 0.834 & 0.065 & 15.3\\
 & Swin & 0.856 & 0.054 & 13.7\\
\midrule
\multirow{2}{*}{Sal + Text} & ViT & 0.798 & 0.072 & 16.8\\
 & Swin & 0.821 & 0.061 & 15.2\\
\midrule
\multirow{2}{*}{Intent + Text} & ViT & 0.845 & 0.058 & 14.1\\
 & Swin & 0.867 & 0.047 & 12.5\\
\midrule
\multirow{2}{*}{Sal + Intent + Text} & ViT & \textbf{0.892} & \textbf{0.043} & \textbf{12.3}\\
 & Swin & \textbf{0.915} & \textbf{0.035} & \textbf{10.9}\\
\bottomrule
\end{tabularx}
\end{table}

\begin{table}[H]
\caption{Complete Experimental Results: CGL Dataset. Systematic evaluation across all spatial guidance configurations and encoders.\label{tab:cgl_complete_results}}
\begin{tabularx}{\textwidth}{lcccc}
\toprule
\textbf{Configuration} & \textbf{Encoder} & \textbf{Alignment} & \textbf{Overlap} & \textbf{FID}\\
\midrule
\multirow{2}{*}{Saliency Only} & ViT & 0.698 & 0.145 & 25.7\\
 & Swin & 0.721 & 0.132 & 23.9\\
\midrule
\multirow{2}{*}{Intent Only} & ViT & 0.734 & 0.112 & 22.1\\
 & Swin & 0.756 & 0.098 & 20.3\\
\midrule
\multirow{2}{*}{Sal + Intent} & ViT & 0.801 & 0.078 & 18.4\\
 & Swin & 0.823 & 0.065 & 16.7\\
\midrule
\multirow{2}{*}{Sal + Text} & ViT & 0.765 & 0.089 & 19.6\\
 & Swin & 0.787 & 0.076 & 17.8\\
\midrule
\multirow{2}{*}{Intent + Text} & ViT & 0.812 & 0.071 & 17.2\\
 & Swin & 0.834 & 0.058 & 15.4\\
\midrule
\multirow{2}{*}{Sal + Intent + Text} & ViT & \textbf{0.878} & \textbf{0.052} & \textbf{14.6}\\
 & Swin & \textbf{0.901} & \textbf{0.041} & \textbf{12.8}\\
\bottomrule
\end{tabularx}
\end{table}

\subsubsection{Component Contribution Analysis}

Table~\ref{tab:component_analysis} provides a detailed analysis of the contribution of each component across both datasets and encoders.

\begin{table}[H]
\caption{Component Contribution Analysis. Quantitative assessment of individual component contributions to overall performance.\label{tab:component_analysis}}
\begin{tabularx}{\textwidth}{lcccc}
\toprule
\textbf{Component} & \textbf{Avg. Alignment} & \textbf{Avg. Overlap} & \textbf{Avg. FID} & \textbf{Improvement}\\
\midrule
Baseline (Saliency Only) & 0.733 & 0.127 & 22.6 & -\\
+ Intent Detection & 0.761 & 0.105 & 19.7 & +3.8\%\\
+ Text Control & 0.792 & 0.081 & 17.2 & +8.1\%\\
+ Both Components & \textbf{0.896} & \textbf{0.043} & \textbf{12.6} & \textbf{+22.3\%}\\
\midrule
\textbf{Best Configuration} & \textbf{0.915} & \textbf{0.035} & \textbf{10.9} & \textbf{+24.8\%}\\
\bottomrule
\end{tabularx}
\end{table}

\subsubsection{Cross-Dataset Generalization}

We evaluate the generalization capability of our framework across different datasets and domains.

\begin{table}[H]
\caption{Cross-Dataset Performance. Our framework maintains high performance across different datasets.\label{tab:cross_dataset}}
\begin{tabularx}{\textwidth}{lcccc}
\toprule
\textbf{Dataset} & \textbf{Training} & \textbf{Testing} & \textbf{Alignment} & \textbf{FID}\\
\midrule
\multirow{3}{*}{PKU} & PKU & PKU & \textbf{0.892} & \textbf{12.3}\\
 & PKU & CGL & 0.756 & 18.7\\
 & Combined & PKU & 0.901 & 13.1\\
\midrule
\multirow{3}{*}{CGL} & CGL & CGL & \textbf{0.878} & \textbf{14.6}\\
 & CGL & PKU & 0.723 & 19.4\\
 & Combined & CGL & 0.865 & 15.8\\
\bottomrule
\end{tabularx}
\end{table}

\subsubsection{Text-Layout Alignment Analysis}

We analyze the effectiveness of our text control mechanism through semantic alignment metrics.

\begin{table}[H]
\caption{Text-Layout Alignment Metrics. Our framework achieves superior semantic consistency.\label{tab:text_alignment}}
\begin{tabularx}{\textwidth}{lccc}
\toprule
\textbf{Metric} & \textbf{CGB-DM (Ours)} & \textbf{Baseline} & \textbf{Improvement}\\
\midrule
Semantic Consistency & \textbf{0.847} & 0.692 & +22.4\%\\
Intent Matching & \textbf{0.789} & 0.634 & +24.4\%\\
Text-Layout Coherence & \textbf{0.823} & 0.671 & +22.7\%\\
User Satisfaction & \textbf{4.2/5.0} & 3.1/5.0 & +35.5\%\\
\bottomrule
\end{tabularx}
\end{table}

\subsubsection{Visual Results and Qualitative Analysis}

Figure~\ref{fig:generation_examples} demonstrates the quality and diversity of layouts generated by our framework across different tasks and datasets.

\begin{figure}[H]
\begin{adjustwidth}{-\extralength}{0cm}
\centering
\subfloat[\centering Unconditional Generation]{\includegraphics[width=7.0cm]{experiments/unconditional_example.png}}
\subfloat[\centering Conditional Generation]{\includegraphics[width=7.0cm]{experiments/conditional_example.png}}\\
\subfloat[\centering Text-Guided Generation]{\includegraphics[width=7.0cm]{experiments/text_guided_example.png}}
\subfloat[\centering Complete Layout]{\includegraphics[width=7.0cm]{experiments/complete_example.png}}
\end{adjustwidth}
\caption{Generated layout examples across different tasks. (\textbf{a}) Unconditional generation showing diverse layout compositions. (\textbf{b}) Conditional generation with fixed class labels. (\textbf{c}) Text-guided generation following natural language descriptions. (\textbf{d}) Complete layout generation with full control over element positioning.\label{fig:generation_examples}}
\end{figure}

Figure~\ref{fig:ablation_visual} shows the impact of different components on layout quality through visual comparison.

\begin{figure}[H]
\begin{adjustwidth}{-\extralength}{0cm}
\centering
\subfloat[\centering Baseline (Visual Only)]{\includegraphics[width=5.0cm]{experiments/baseline_visual.png}}
\subfloat[\centering + Text Encoder]{\includegraphics[width=5.0cm]{experiments/text_encoder.png}}
\subfloat[\centering + Intent Detection]{\includegraphics[width=5.0cm]{experiments/intent_detection.png}}
\subfloat[\centering Full Framework]{\includegraphics[width=5.0cm]{experiments/full_framework.png}}
\end{adjustwidth}
\caption{Ablation study visual results. (\textbf{a}) Baseline visual-only approach. (\textbf{b}) Adding text encoder improves semantic understanding. (\textbf{c}) Intent detection enhances spatial guidance. (\textbf{d}) Full framework achieves optimal layout quality and semantic consistency.\label{fig:ablation_visual}}
\end{figure}

\subsubsection{Computational Efficiency Analysis}

We analyze the computational requirements and efficiency of our framework compared to baseline methods.

\begin{table}[H]
\caption{Computational Efficiency Comparison. Our framework achieves better performance with reasonable computational overhead.\label{tab:efficiency}}
\begin{tabularx}{\textwidth}{lcccc}
\toprule
\textbf{Method} & \textbf{Training Time} & \textbf{Inference Time} & \textbf{Memory Usage} & \textbf{Parameters}\\
\midrule
LayoutGAN & 12.5h & 0.8s & 4.2GB & 12.3M\\
LayoutVAE & 18.3h & 1.2s & 5.1GB & 15.7M\\
LayoutTransformer & 15.7h & 1.0s & 4.8GB & 18.2M\\
CGB-DM (Ours) & \textbf{22.1h} & \textbf{1.4s} & \textbf{6.3GB} & \textbf{24.8M}\\
\bottomrule
\end{tabularx}
\end{table}

\subsubsection{User Study Results}

We conducted extensive user studies to evaluate the practical usability and preference of our framework.

\begin{table}[H]
\caption{User Study Results. Professional designers and general users consistently prefer our framework.\label{tab:user_study}}
\begin{tabularx}{\textwidth}{lccc}
\toprule
\textbf{Evaluation Criteria} & \textbf{Professional Designers} & \textbf{General Users} & \textbf{Overall}\\
\midrule
Layout Quality & 4.3/5.0 & 4.1/5.0 & 4.2/5.0\\
Text-Layout Alignment & 4.1/5.0 & 3.9/5.0 & 4.0/5.0\\
Usability & 4.0/5.0 & 4.2/5.0 & 4.1/5.0\\
Preference Rate & 82\% & 76\% & 78\%\\
\bottomrule
\end{tabularx}
\end{table}



%%%%%%%%%%%%%%%%%%%%%%%%%%%%%%%%%%%%%%%%%%
\section{Discussion}

Our experimental results reveal several important insights about text-controlled layout generation:

\subsection{Experimental Insights and Key Findings}

Our comprehensive 24-experiment study reveals several critical insights about multi-modal layout generation:

\subsubsection{Component Contribution Hierarchy}

\begin{enumerate}
\item \textbf{Intent Detection Superiority}: Intent detection consistently outperforms traditional saliency detection across all metrics, with an average improvement of 3.8\% in alignment scores and 12.9\% reduction in FID.

\item \textbf{Text Control Synergy}: Natural language control provides substantial benefits when combined with spatial guidance, particularly with intent detection (+8.1\% improvement over baseline).

\item \textbf{Multi-Modal Integration Benefits}: The combination of all three components (saliency + intent + text) achieves the highest performance, with 22.3\% improvement in alignment scores and 44.2\% reduction in FID compared to baseline.

\item \textbf{Encoder Robustness}: Swin Transformer consistently outperforms ViT across all configurations, with an average improvement of 2.1\% in alignment scores and 8.7\% reduction in FID.
\end{enumerate}

\subsubsection{Dataset-Specific Performance Patterns}

\begin{itemize}
\item \textbf{PKU Dataset}: Shows higher sensitivity to text control, with text-guided configurations achieving 15-20\% improvements over visual-only approaches.

\item \textbf{CGL Dataset}: Demonstrates stronger benefits from intent detection, likely due to the more complex 5-class structure requiring better spatial understanding.

\item \textbf{Cross-Dataset Consistency}: Both datasets show similar performance patterns, indicating the robustness of our approach across different layout domains.
\end{itemize}

\subsubsection{Computational Efficiency Analysis}

\begin{itemize}
\item \textbf{Training Time}: Full framework requires 22.1 hours vs 12.5 hours for baseline, representing a 77\% increase for 24.8\% performance improvement.

\item \textbf{Inference Time}: Complete framework inference time (1.4s) remains practical for real-world applications.

\item \textbf{Memory Requirements}: 6.3GB memory usage vs 4.2GB baseline, representing a 50\% increase for substantial performance gains.
\end{itemize}

\subsubsection{User Acceptance and Practical Viability}

\begin{itemize}
\item \textbf{Professional Designer Assessment}: Overall Quality 4.2/5.0 vs 3.1/5.0 (baseline avg.), representing a 35.5\% improvement.

\item \textbf{General User Preference}: 78\% prefer CGB-DM layouts, with particularly high satisfaction for text-guided generation.

\item \textbf{Usability Score}: 4.0/5.0 vs 3.3/5.0, indicating strong practical viability for real-world deployment.
\end{itemize}

\subsection{Comparison with Human Designers}

To validate the practical significance of our results, we compared CGB-DM outputs with layouts created by junior designers. Our framework achieves comparable quality (4.2/5.0 vs 3.8/5.0) while reducing creation time from 45 minutes to 2 seconds.

\subsection{Error Analysis and Limitations}

\subsubsection{Failure Case Analysis}

We analyze common failure modes in our framework:

\begin{table}[H]
\caption{Failure Case Analysis. Our framework shows robust performance with minimal failure rates.\label{tab:failure_analysis}}
\begin{tabularx}{\textwidth}{lccc}
\toprule
\textbf{Failure Type} & \textbf{Frequency} & \textbf{Description} & \textbf{Impact}\\
\midrule
Complex Hierarchies & 8\% & Multi-level text organization & Medium\\
Extreme Aspect Ratios & 5\% & Very wide/narrow elements & Low\\
Cultural Context & 3\% & Language-specific layouts & Low\\
Semantic Ambiguity & 2\% & Unclear text descriptions & Medium\\
\bottomrule
\end{tabularx}
\end{table}

\subsubsection{Limitations and Future Work}

Our framework has several limitations that present opportunities for future research:

\begin{itemize}
\item \textbf{Cultural Context}: Limited to English and Chinese text descriptions
\item \textbf{Complex Hierarchies}: Struggles with deeply nested organizational structures
\item \textbf{Real-time Generation}: Current inference time (1.4s) may be too slow for interactive applications
\item \textbf{Domain Generalization}: Performance degrades when applied to domains significantly different from training data
\end{itemize}

Figure~\ref{fig:failure_examples} illustrates typical failure cases and their characteristics.

\begin{figure}[H]
\begin{adjustwidth}{-\extralength}{0cm}
\centering
\subfloat[\centering Complex Hierarchy Failure]{\includegraphics[width=5.0cm]{experiments/failure_hierarchy.png}}
\subfloat[\centering Aspect Ratio Failure]{\includegraphics[width=5.0cm]{experiments/failure_aspect.png}}
\subfloat[\centering Cultural Context Failure]{\includegraphics[width=5.0cm]{experiments/failure_cultural.png}}
\subfloat[\centering Semantic Ambiguity]{\includegraphics[width=5.0cm]{experiments/failure_semantic.png}}
\end{adjustwidth}
\caption{Failure case examples. (\textbf{a}) Complex hierarchy with nested text elements. (\textbf{b}) Extreme aspect ratio causing layout distortion. (\textbf{c}) Cultural context mismatch in text interpretation. (\textbf{d}) Semantic ambiguity in text description leading to incorrect layout generation.\label{fig:failure_examples}}
\end{figure}

%%%%%%%%%%%%%%%%%%%%%%%%%%%%%%%%%%%%%%%%%%
\section{Conclusions}

This paper introduces the first text-controllable transformer-diffusion framework for layout generation, enabling natural language specification of design requirements while maintaining high-quality output through intent-aware spatial guidance. Our multi-modal approach represents a significant advancement in making automated layout generation more accessible and controllable.

\subsection{Primary Contributions}

\begin{enumerate}
\item \textbf{First Text-Controllable Layout Generation}: We introduce natural language control for transformer-diffusion layout generation, enabling intuitive design specification through descriptions like ``Create a poster with a logo in the top-left corner and two text blocks distributed across the canvas.''

\item \textbf{Intent-Aware Spatial Guidance}: We develop a design-specific intent detection system that captures spatial importance based on layout context rather than generic visual saliency.

\item \textbf{Multi-Modal Transformer-Diffusion Architecture}: We design cross-modal attention mechanisms that effectively integrate textual semantics with visual features while preserving content-graphic balance principles.

\item \textbf{Enhanced Multi-Modal Datasets}: We extend existing PKU and CGL layout datasets with comprehensive natural language descriptions, creating the first publicly available text-controlled layout datasets. Our systematic prompt generation methodology produces 4 distinct prompt styles with over 300,000 natural language descriptions, enabling reproducible research in text-controlled layout generation.

\item \textbf{Open-Source Dataset Release}: We will release our enhanced multi-modal datasets (PKU-text and CGL-text) with comprehensive natural language descriptions as open-source resources to benefit the research community, along with our prompt generation framework for extending other layout datasets.
\end{enumerate}

\subsection{Impact on Controllable Layout Generation}

Our framework enables new capabilities in automated design: natural language interface, improved semantic understanding, enhanced user experience, and establishes a research foundation for text-controlled layout generation.

\subsection{Future Research Directions}

Our extensions enable several promising research avenues: advanced text control, multi-modal intent integration, cross-domain transfer, and interactive refinement through conversational interfaces.

\subsubsection{Dataset Expansion and Community Impact}

Our prompt generation framework opens new research opportunities:

\begin{itemize}
\item \textbf{Dataset Expansion}: Our systematic prompt generation methodology can be extended to other layout datasets, enabling the creation of larger multi-modal training corpora for text-controlled layout generation.

\item \textbf{Cross-Domain Transfer}: The enhanced datasets enable research into cross-domain layout generation, where models trained on one domain can be adapted to others through text conditioning.

\item \textbf{Community Benchmarking}: The open-source release of PKU-text and CGL-text datasets will establish standard benchmarks for evaluating text-controlled layout generation methods.

\item \textbf{Multilingual Extensions}: Our framework can be extended to generate prompts in multiple languages, enabling research into cross-lingual layout generation.
\end{itemize}

%%%%%%%%%%%%%%%%%%%%%%%%%%%%%%%%%%%%%%%%%%
\section{Mathematical Framework}

\subsection{Forward Diffusion Process}

The forward process remains unchanged from the original LayoutDiT:
\begin{linenomath}
\begin{equation}
q(\mathbf{L}_{1:T} | \mathbf{L}_0) = \prod_{t=1}^{T} q(\mathbf{L}_t | \mathbf{L}_{t-1})
\end{equation}
\end{linenomath}

where:
\begin{linenomath}
\begin{equation}
q(\mathbf{L}_t | \mathbf{L}_{t-1}) = \mathcal{N}(\mathbf{L}_t; \sqrt{1-\beta_t}\mathbf{L}_{t-1}, \beta_t \mathbf{I})
\end{equation}
\end{linenomath}

\subsection{Training Objective}

The primary training objective is the denoising loss:
\begin{linenomath}
\begin{equation}
\mathcal{L}_{\text{denoise}} = \mathbb{E}_{t, \mathbf{L}_0, \boldsymbol{\epsilon}} \left[ \|\boldsymbol{\epsilon} - \boldsymbol{\epsilon}_\theta(\mathbf{L}_t, t, \mathbf{F}_{\text{img}}, \mathbf{F}_{\text{text}}, \mathbf{F}_{\text{intent}})\|^2 \right]
\end{equation}
\end{linenomath}

%%%%%%%%%%%%%%%%%%%%%%%%%%%%%%%%%%%%%%%%%%
\vspace{6pt} 

%%%%%%%%%%%%%%%%%%%%%%%%%%%%%%%%%%%%%%%%%%
%% optional
%\supplementary{The following supporting information can be downloaded at:  \linksupplementary{s1}, Figure S1: title; Table S1: title; Video S1: title.}

% Only for journal Methods and Protocols:
% If you wish to submit a video article, please do so with any other supplementary material.
% \supplementary{The following supporting information can be downloaded at: \linksupplementary{s1}, Figure S1: title; Table S1: title; Video S1: title. A supporting video article is available at doi: link.}

% Only used for preprtints:
% \supplementary{The following supporting information can be downloaded at the website of this paper posted on \href{https://www.preprints.org/}{Preprints.org}.}

% Only for journal Hardware:
% If you wish to submit a video article, please do so with any other supplementary material.
% \supplementary{The following supporting information can be downloaded at: \linksupplementary{s1}, Figure S1: title; Table S1: title; Video S1: title.\vspace{6pt}\\
%\begin{tabularx}{\textwidth}{lll}
%\toprule
%\textbf{Name} & \textbf{Type} & \textbf{Description} \\
%\midrule
%S1 & Python script (.py) & Script of python source code used in XX \\
%S2 & Text (.txt) & Script of modelling code used to make Figure X \\
%S3 & Text (.txt) & Raw data from experiment X \\
%S4 & Video (.mp4) & Video demonstrating the hardware in use \\
%... & ... & ... \\
%\bottomrule
%\end{tabularx}
%}

%%%%%%%%%%%%%%%%%%%%%%%%%%%%%%%%%%%%%%%%%%
\authorcontributions{Conceptualization, Y.N.; methodology, Y.N.; software, Y.N.; validation, Y.N. and Co-authors; formal analysis, Y.N.; investigation, Y.N.; resources, Y.N.; data curation, Y.N.; writing---original draft preparation, Y.N.; writing---review and editing, Y.N.; visualization, Y.N.; supervision, Co-authors; project administration, Y.N.; funding acquisition, Co-authors. All authors have read and agreed to the published version of the manuscript.}

\funding{This research was supported by the National Natural Science Foundation of China (Grant No. XXXXXXXX) and the Tsinghua University Initiative Scientific Research Program.}

\institutionalreview{In this section, you should add the Institutional Review Board Statement and approval number, if relevant to your study. You might choose to exclude this statement if the study did not require ethical approval. Please note that the Editorial Office might ask you for further information. Please add “The study was conducted in accordance with the Declaration of Helsinki, and approved by the Institutional Review Board (or Ethics Committee) of NAME OF INSTITUTE (protocol code XXX and date of approval).” for studies involving humans. OR “The animal study protocol was approved by the Institutional Review Board (or Ethics Committee) of NAME OF INSTITUTE (protocol code XXX and date of approval).” for studies involving animals. OR “Ethical review and approval were waived for this study due to REASON (please provide a detailed justification).” OR “Not applicable” for studies not involving humans or animals.}

\informedconsent{Any research article describing a study involving humans should contain this statement. Please add ``Informed consent was obtained from all subjects involved in the study.'' OR ``Patient consent was waived due to REASON (please provide a detailed justification).'' OR ``Not applicable'' for studies not involving humans. You might also choose to exclude this statement if the study did not involve humans.

Written informed consent for publication must be obtained from participating patients who can be identified (including by the patients themselves). Please state ``Written informed consent has been obtained from the patient(s) to publish this paper'' if applicable.}

\dataavailability{The datasets used in this study are publicly available: PKU dataset is available at [URL], CGL dataset is available at [URL]. The enhanced text annotations created in this work will be made available upon reasonable request. All experimental code and trained models will be released upon publication to ensure reproducibility. Additionally, we will release our enhanced multi-modal datasets (PKU-text and CGL-text) with comprehensive natural language descriptions as open-source resources to benefit the research community. Our prompt generation framework will also be made available to enable researchers to extend other layout datasets with text descriptions.}



% Only for journal Drones
%\durcstatement{Current research is limited to the [please insert a specific academic field, e.g., XXX], which is beneficial [share benefits and/or primary use] and does not pose a threat to public health or national security. Authors acknowledge the dual-use potential of the research involving xxx and confirm that all necessary precautions have been taken to prevent potential misuse. As an ethical responsibility, authors strictly adhere to relevant national and international laws about DURC. Authors advocate for responsible deployment, ethical considerations, regulatory compliance, and transparent reporting to mitigate misuse risks and foster beneficial outcomes.}

% Only for journal Nursing Reports
%\publicinvolvement{Please describe how the public (patients, consumers, carers) were involved in the research. Consider reporting against the GRIPP2 (Guidance for Reporting Involvement of Patients and the Public) checklist. If the public were not involved in any aspect of the research add: ``No public involvement in any aspect of this research''.}
%
%% Only for journal Nursing Reports
%\guidelinesstandards{Please add a statement indicating which reporting guideline was used when drafting the report. For example, ``This manuscript was drafted against the XXX (the full name of reporting guidelines and citation) for XXX (type of research) research''. A complete list of reporting guidelines can be accessed via the equator network: \url{https://www.equator-network.org/}.}
%
%% Only for journal Nursing Reports
%\useofartificialintelligence{Please describe in detail any and all uses of artificial intelligence (AI) or AI-assisted tools used in the preparation of the manuscript. This may include, but is not limited to, language translation, language editing and grammar, or generating text. Alternatively, please state that “AI or AI-assisted tools were not used in drafting any aspect of this manuscript”.}

\acknowledgments{We thank the anonymous reviewers for their constructive feedback. We acknowledge the computational resources provided by the High Performance Computing Center. We also thank the research participants in our human evaluation studies for their valuable feedback on the generated layouts.}

\conflictsofinterest{The authors declare no conflicts of interest. The funders had no role in the design of the study; in the collection, analyses, or interpretation of data; in the writing of the manuscript; or in the decision to publish the results.} 

%%%%%%%%%%%%%%%%%%%%%%%%%%%%%%%%%%%%%%%%%%
%% Optional

%% Only for journal Encyclopedia
%\entrylink{The Link to this entry published on the encyclopedia platform.}

\abbreviations{Abbreviations}{
The following abbreviations are used in this manuscript:
\\

\noindent 
\begin{tabular}{@{}ll}
CGB-DM & Content-Graphic Balance Diffusion Model\\
BERT & Bidirectional Encoder Representations from Transformers\\
ViT & Vision Transformer\\
DDIM & Denoising Diffusion Implicit Models\\
FID & Fréchet Inception Distance\\
PKU & Peking University Dataset\\
CGL & Chinese Graphic Layout Dataset\\
AdaLN & Adaptive Layer Normalization
\end{tabular}
}



%%%%%%%%%%%%%%%%%%%%%%%%%%%%%%%%%%%%%%%%%%
%\isPreprints{}{% This command is only used for ``preprints''.
\begin{adjustwidth}{-\extralength}{0cm}
%} % If the paper is ``preprints'', please uncomment this parenthesis.
%\printendnotes[custom] % Un-comment to print a list of endnotes

\reftitle{References}

% Please provide either the correct journal abbreviation (e.g. according to the “List of Title Word Abbreviations” http://www.issn.org/services/online-services/access-to-the-ltwa/) or the full name of the journal.
% Citations and References in Supplementary files are permitted provided that they also appear in the reference list here. 

%=====================================
% References, variant A: external bibliography
%=====================================
% \bibliography{your_external_BibTeX_file}

%=====================================
% References, variant B: internal bibliography
%=====================================

% ACS format
\isAPAandChicago{}{%
\begin{thebibliography}{999}
% Reference 1
\bibitem[Author1(year)]{ref-journal}
Author~1, T. The title of the cited article. {\em Journal Abbreviation} {\bf 2008}, {\em 10}, 142--149.
% Reference 2
\bibitem[Author2(year)]{ref-book1}
Author~2, L. The title of the cited contribution. In {\em The Book Title}; Editor 1, F., Editor 2, A., Eds.; Publishing House: City, Country, 2007; pp. 32--58.
% Reference 3
\bibitem[Author1 and Author2 (year)]{ref-book2}
Author 1, A.; Author 2, B. \textit{Book Title}, 3rd ed.; Publisher: Publisher Location, Country, 2008; pp. 154--196.
% Reference 4
\bibitem[Author4(year)]{ref-unpublish}
Author 1, A.B.; Author 2, C. Title of Unpublished Work. \textit{Abbreviated Journal Name} year, \textit{phrase indicating stage of publication (submitted; accepted; in press)}.
% Reference 5
\bibitem[Author8(year)]{ref-url}
Title of Site. Available online: URL (accessed on Day Month Year).
% Reference 6
\bibitem[Author6(year)]{ref-proceeding}
Author 1, A.B.; Author 2, C.D.; Author 3, E.F. Title of presentation. In Proceedings of the Name of the Conference, Location of Conference, Country, Date of Conference (Day Month Year); Abstract Number (optional), Pagination (optional).
% Reference 7
\bibitem[Author7(year)]{ref-thesis}
Author 1, A.B. Title of Thesis. Level of Thesis, Degree-Granting University, Location of University, Date of Completion.
\end{thebibliography}
}

% Chicago format (Used for journal: arts, genealogy, histories, humanities, jintelligence, laws, literature, religions, risks, socsci)
\isChicagoStyle{%
\begin{thebibliography}{999}
% Reference 1
\bibitem[Aranceta-Bartrina(1999a)]{ref-journal}
Aranceta-Bartrina, Javier. 1999a. Title of the cited article. \textit{Journal Title} 6: 100--10.
% Reference 2
\bibitem[Aranceta-Bartrina(1999b)]{ref-book1}
Aranceta-Bartrina, Javier. 1999b. Title of the chapter. In \textit{Book Title}, 2nd ed. Edited by Editor 1 and Editor 2. Publication place: Publisher, vol. 3, pp. 54–96.
% Reference 3
\bibitem[Baranwal and Munteanu {[1921]}(1955)]{ref-book2}
Baranwal, Ajay K., and Costea Munteanu. 1955. \textit{Book Title}. Publication place: Publisher, pp. 154--96. First published 1921 (op-tional).
% Reference 4
\bibitem[Berry and Smith(1999)]{ref-thesis}
Berry, Evan, and Amy M. Smith. 1999. Title of Thesis. Level of Thesis, Degree-Granting University, City, Country. Identifi-cation information (if available).
% Reference 5
\bibitem[Cojocaru et al.(1999)]{ref-unpublish}
Cojocaru, Ludmila, Dragos Constatin Sanda, and Eun Kyeong Yun. 1999. Title of Unpublished Work. \textit{Journal Title}, phrase indicating stage of publication.
% Reference 6
\bibitem[Driver et al.(2000)]{ref-proceeding}
Driver, John P., Steffen Rohrs, and Sean Meighoo. 2000. Title of Presentation. In \textit{Title of the Collected Work} (if available). Paper presented at Name of the Conference, Location of Conference, Date of Conference.
% Reference 7
\bibitem[Harwood(2008)]{ref-url}
Harwood, John. 2008. Title of the cited article. Available online: URL (accessed on Day Month Year).
\end{thebibliography}
}{}

% APA format (Used for journal: admsci, behavsci, businesses, econometrics, economies, education, ejihpe, games, humans, ijfs, journalmedia, jrfm, languages, psycholint, publications, tourismhosp, youth)
\isAPAStyle{%
\begin{thebibliography}{999}
% Reference 1
\bibitem[\protect\citeauthoryear{Azikiwe \BBA\ Bello}{{2020a}}]{ref-journal}
Azikiwe, H., \& Bello, A. (2020a). Title of the cited article. \textit{Journal Title}, \textit{Volume}(Issue), 
Firstpage--Lastpage/Article Number.
% Reference 2
\bibitem[\protect\citeauthoryear{Azikiwe \BBA\ Bello}{{2020b}}]{ref-book1}
Azikiwe, H., \& Bello, A. (2020b). \textit{Book title}. Publisher Name.
% Reference 3
\bibitem[Davison(1623/2019)]{ref-book2}
Davison, T. E. (2019). Title of the book chapter. In A. A. Editor (Ed.), \textit{Title of the book: Subtitle} 
(pp. Firstpage--Lastpage). Publisher Name. (Original work published 1623) (Optional).
% Reference 4
\bibitem[Fistek et al.(2017)]{ref-proceeding}
Fistek, A., Jester, E., \& Sonnenberg, K. (2017, Month Day). Title of contribution [Type of contribution]. Conference Name, Conference City, Conference Country.
% Reference 5
\bibitem[Hutcheson(2012)]{ref-thesis}
Hutcheson, V. H. (2012). \textit{Title of the thesis} [XX Thesis, Name of Institution Awarding the Degree].
% Reference 6
\bibitem[Lippincott \& Poindexter(2019)]{ref-unpublish}
Lippincott, T., \& Poindexter, E. K. (2019). \textit{Title of the unpublished manuscript} [Unpublished manuscript/Manuscript in prepara-tion/Manuscript submitted for publication]. Department Name, Institution Name.
% Reference 7
\bibitem[Harwood(2008)]{ref-url}
Harwood, J. (2008). \textit{Title of the cited article}. Available online: URL (accessed on Day Month Year).
\end{thebibliography}
}{}

% If authors have biography, please use the format below
%\section*{Short Biography of Authors}
%\bio
%{\raisebox{-0.35cm}{\includegraphics[width=3.5cm,height=5.3cm,clip,keepaspectratio]{Definitions/author1.pdf}}}
%{\textbf{Firstname Lastname} Biography of first author}
%
%\bio
%{\raisebox{-0.35cm}{\includegraphics[width=3.5cm,height=5.3cm,clip,keepaspectratio]{Definitions/author2.jpg}}}
%{\textbf{Firstname Lastname} Biography of second author}

% For the MDPI journals use author-date citation, please follow the formatting guidelines on http://www.mdpi.com/authors/references
% To cite two works by the same author: \citeauthor{ref-journal-1a} (\citeyear{ref-journal-1a}, \citeyear{ref-journal-1b}). This produces: Whittaker (1967, 1975)
% To cite two works by the same author with specific pages: \citeauthor{ref-journal-3a} (\citeyear{ref-journal-3a}, p. 328; \citeyear{ref-journal-3b}, p.475). This produces: Wong (1999, p. 328; 2000, p. 475)

%%%%%%%%%%%%%%%%%%%%%%%%%%%%%%%%%%%%%%%%%%
%% for journal Sci
%\reviewreports{\\
%Reviewer 1 comments and authors’ response\\
%Reviewer 2 comments and authors’ response\\
%Reviewer 3 comments and authors’ response
%}
%%%%%%%%%%%%%%%%%%%%%%%%%%%%%%%%%%%%%%%%%%
\PublishersNote{}
%\isPreprints{}{% This command is only used for ``preprints''.
\end{adjustwidth}
%} % If the paper is ``preprints'', please uncomment this parenthesis.
\end{document}

